% GENERATED FILE. Edit canonical lessons, then run npm run generate:books.
% canonical-source-hash: fnv1a64:38ff9f1c77611ea2
% canonical-lessons: GE-C293-brennen, GE-C293-rauch, GE-C293-flamme, GE-C293-loeschen, GE-C293-waldbrand, GE-R293-review-286-293

\chapter{Es brennt}
\label{ch:ge-es-brennt}
\hlchaptermodality{293}

\begin{chapteropening}
\textbf{By the end of this chapter:} \emph{I can report a fire: to burn, smoke, a flame, to put out and a forest fire.}

Complete the last lesson of chapter 293: I can report a fire: to burn, smoke, a flame, to put out and a forest fire.
\end{chapteropening}

\section[brennen]{brennen — to burn}

\label{lesson:GE-C293-brennen}



\begin{quote}
Before the new one: say the German for to evacuate, then the German for to survive.
\end{quote}

\subsection*{You'll want to know: brennen}
\textbf{brennen} — \textquotedblleft{}to burn\textquotedblright{}.

\begin{quote}
Das Haus brannte die ganze Nacht.
\end{quote}

\begin{quote}
\emph{The house burned all night.}
\end{quote}

Mixed verb: \emph{brannte}, \emph{hat gebrannt}. \emph{Es brennt!} — fire!

\begin{cousinweb}
\emph{brennen} is an inherited Germanic verb: EtymDB traces Old High German \emph{brennen} to Proto-Germanic \emph{brannijaną}, built on \emph{brinnaną} \textquotedblleft{}to burn\textquotedblright{}, and lists Old English \emph{birnan} \textquotedblleft{}to burn\textquotedblright{} among its cousins: the family of English \emph{burn}.
\end{cousinweb}

\begin{grammarlens}[title={What you've built}]
One more word for telling, explaining and describing at B1.
\end{grammarlens}

\subsection*{Your turn}
\begin{itemize}
\raggedright
  \item \emph{Say it:} \emph{brennen}
  \item \emph{Say it:} \emph{brennen}, once more
  \item \emph{Say it:} a sentence of your own with \emph{brennen}
  \item \emph{Recall:} say the German for to evacuate, then the German for to survive, then say \emph{brennen} again
\end{itemize}

\begin{tcolorbox}[breakable,colback=teal!4,colframe=teal!35!black,title={Before you move on}]
What does brennen mean? (\textquotedblleft{}To burn\textquotedblright{}.) Say it once more.
\end{tcolorbox}
\section[der Rauch]{der Rauch — the smoke}

\label{lesson:GE-C293-rauch}



\begin{quote}
Before the new one: say the German for to survive, then the German for to burn.
\end{quote}

\subsection*{You'll want to know: der Rauch}
\textbf{der Rauch} — \textquotedblleft{}the smoke\textquotedblright{}.

\begin{quote}
Aus dem Fenster kam dicker Rauch.
\end{quote}

\begin{quote}
\emph{Thick smoke came out of the window.}
\end{quote}

No plural. \emph{Rauchen} — to smoke — comes from it.

\begin{cousinweb}
\emph{Rauch} is an inherited Germanic word; EtymDB lists English \emph{reek} and Dutch \emph{rook} \textquotedblleft{}smoke\textquotedblright{} among its cousins.
\end{cousinweb}

\begin{grammarlens}[title={What you've built}]
One more word for telling, explaining and describing at B1.
\end{grammarlens}

\subsection*{Your turn}
\begin{itemize}
\raggedright
  \item \emph{Say it:} \emph{der Rauch}
  \item \emph{Say it:} \emph{der Rauch}, once more
  \item \emph{Say it:} a sentence of your own with \emph{der Rauch}
  \item \emph{Recall:} say the German for to survive, then the German for to burn, then say \emph{der Rauch} again
\end{itemize}

\begin{tcolorbox}[breakable,colback=teal!4,colframe=teal!35!black,title={Before you move on}]
What does der Rauch mean? (\textquotedblleft{}The smoke\textquotedblright{}.) Say it once more.
\end{tcolorbox}
\section[die Flamme]{die Flamme — the flame}

\label{lesson:GE-C293-flamme}



\begin{quote}
Before the new one: say the German for to burn, then the German for the smoke.
\end{quote}

\subsection*{You'll want to know: die Flamme}
\textbf{die Flamme} — \textquotedblleft{}the flame\textquotedblright{}.

\begin{quote}
Die Flammen waren schon hoch über dem Dach.
\end{quote}

\begin{quote}
\emph{The flames were already high above the roof.}
\end{quote}

Plural: \emph{die Flammen}. \emph{In Flammen stehen} — to be on fire.

\begin{cousinweb}
\emph{Flamme} was borrowed from Latin \emph{flamma} \textquotedblleft{}a flame, a blaze\textquotedblright{} (Lewis \& Short); EtymDB records the loan already in Old High German \emph{flamma}.
\end{cousinweb}

\begin{grammarlens}[title={What you've built}]
One more word for telling, explaining and describing at B1.
\end{grammarlens}

\subsection*{Your turn}
\begin{itemize}
\raggedright
  \item \emph{Say it:} \emph{die Flamme}
  \item \emph{Say it:} \emph{die Flamme}, once more
  \item \emph{Say it:} a sentence of your own with \emph{die Flamme}
  \item \emph{Recall:} say the German for to burn, then the German for the smoke, then say \emph{die Flamme} again
\end{itemize}

\begin{tcolorbox}[breakable,colback=teal!4,colframe=teal!35!black,title={Before you move on}]
What does die Flamme mean? (\textquotedblleft{}The flame\textquotedblright{}.) Say it once more.
\end{tcolorbox}
\section[löschen]{löschen — to put out, to extinguish; to delete}

\label{lesson:GE-C293-loeschen}



\begin{quote}
Before the new one: say the German for the smoke, then the German for the flame.
\end{quote}

\subsection*{You'll want to know: löschen}
\textbf{löschen} — \textquotedblleft{}to put out, to extinguish; to delete\textquotedblright{}.

\begin{quote}
Die Feuerwehr löschte das Feuer in einer Stunde.
\end{quote}

\begin{quote}
\emph{The fire brigade put out the fire in an hour.}
\end{quote}

On a computer it means to delete: \emph{ein Foto löschen}.

\begin{cousinweb}
\emph{löschen} is inherited from Middle High German \emph{leschen}, Old High German \emph{lescan} (EtymDB). The old verb is usually glossed \textquotedblleft{}to go out, to die down\textquotedblright{}; later it came to mean \textquotedblleft{}to make go out\textquotedblright{}.
\end{cousinweb}

\begin{grammarlens}[title={What you've built}]
One more word for telling, explaining and describing at B1.
\end{grammarlens}

\subsection*{Your turn}
\begin{itemize}
\raggedright
  \item \emph{Say it:} \emph{löschen}
  \item \emph{Say it:} \emph{löschen}, once more
  \item \emph{Say it:} a sentence of your own with \emph{löschen}
  \item \emph{Recall:} say the German for the smoke, then the German for the flame, then say \emph{löschen} again
\end{itemize}

\begin{tcolorbox}[breakable,colback=teal!4,colframe=teal!35!black,title={Before you move on}]
What does löschen mean? (\textquotedblleft{}To put out, to extinguish; to delete\textquotedblright{}.) Say it once more.
\end{tcolorbox}
\section[der Waldbrand]{der Waldbrand — the forest fire}

\label{lesson:GE-C293-waldbrand}



\begin{quote}
Before the new one: say the German for the flame, then the German for to put out.
\end{quote}

\subsection*{You'll want to know: der Waldbrand}
\textbf{der Waldbrand} — \textquotedblleft{}the forest fire\textquotedblright{}.

\begin{quote}
Im Sommer gab es viele Waldbrände.
\end{quote}

\begin{quote}
\emph{In summer there were many forest fires.}
\end{quote}

Plural: \emph{die Waldbrände}. \emph{Der Brand} alone is a large fire.

\begin{cousinweb}
\emph{Waldbrand} joins \emph{Wald} \textquotedblleft{}forest\textquotedblright{}, Old High German \emph{wald} (EtymDB), and \emph{Brand} \textquotedblleft{}fire\textquotedblright{}, built on \emph{brennen} \textquotedblleft{}to burn\textquotedblright{}.
\end{cousinweb}

\begin{grammarlens}[title={What you've built}]
One more word for telling, explaining and describing at B1.
\end{grammarlens}

\subsection*{Your turn}
\begin{itemize}
\raggedright
  \item \emph{Say it:} \emph{der Waldbrand}
  \item \emph{Say it:} \emph{der Waldbrand}, once more
  \item \emph{Say it:} a sentence of your own with \emph{der Waldbrand}
  \item \emph{Recall:} say the German for the flame, then the German for to put out, then say \emph{der Waldbrand} again
\end{itemize}

\begin{tcolorbox}[breakable,colback=teal!4,colframe=teal!35!black,title={Before you move on}]
What does der Waldbrand mean? (\textquotedblleft{}The forest fire\textquotedblright{}.) Say it once more.
\end{tcolorbox}
\section[(dialogue)]{Words from chapters 286–293 — a review}

\label{lesson:GE-R293-review-286-293}



\begin{quote}
Say the German for the verdict and the judge.
\end{quote}

\subsection*{Your turn}
Say these aloud:

\begin{itemize}
\raggedright
  \item the verdict, the judge, to sentence, guilty, innocent
  \item to confess, the confession, the prison, to disappear, to flee
  \item to collide, the collision, to crash, to have an accident, the damage
  \item to slip, to trip, to fall, the fall, to bleed
  \item unconscious, faint, consciousness, to choke, to drown
  \item the emergency call, the emergency doctor, the helicopter, the paramedic, the emergency room
  \item to warn, the rescue, to bring out alive, to evacuate, to survive
  \item to burn, the smoke, the flame, to put out, the forest fire
\end{itemize}

\begin{tcolorbox}[breakable,colback=teal!4,colframe=teal!35!black,title={Before you move on}]
The verdict? (\textbf{das Urteil}.) The judge? (\textbf{der Richter}.) To sentence? (\textbf{verurteilen}.) Guilty? (\textbf{schuldig}.) Innocent? (\textbf{unschuldig}.) To confess? (\textbf{gestehen}.) The confession? (\textbf{das Geständnis}.) The prison? (\textbf{das Gefängnis}.) To disappear? (\textbf{verschwinden}.) To flee? (\textbf{fliehen}.) To collide? (\textbf{zusammenstoßen}.) The collision? (\textbf{der Zusammenstoß}.) To crash? (\textbf{prallen}.) To have an accident? (\textbf{verunglücken}.) The damage? (\textbf{der Schaden}.) To slip? (\textbf{ausrutschen}.) To trip? (\textbf{stolpern}.) To fall? (\textbf{stürzen}.) The fall? (\textbf{der Sturz}.) To bleed? (\textbf{bluten}.) Unconscious? (\textbf{bewusstlos}.) Faint? (\textbf{ohnmächtig}.) Consciousness? (\textbf{das Bewusstsein}.) To choke? (\textbf{ersticken}.) To drown? (\textbf{ertrinken}.) The emergency call? (\textbf{der Notruf}.) The emergency doctor? (\textbf{der Notarzt}.) The helicopter? (\textbf{der Hubschrauber}.) The paramedic? (\textbf{der Sanitäter}.) The emergency room? (\textbf{die Notaufnahme}.) To warn? (\textbf{warnen}.) The rescue? (\textbf{die Rettung}.) To bring out alive? (\textbf{bergen}.) To evacuate? (\textbf{evakuieren}.) To survive? (\textbf{überleben}.) To burn? (\textbf{brennen}.) The smoke? (\textbf{der Rauch}.) The flame? (\textbf{die Flamme}.) To put out? (\textbf{löschen}.) The forest fire? (\textbf{der Waldbrand}.)
\end{tcolorbox}
